\subsection{Preliminaries}

Before reviewing prior work it helps to fix a few ideas that recur throughout
the thesis. A convolutional neural network builds up a representation of an image
through stacked local filters, which makes it very good at capturing texture and
short-range structure but, by construction, slow to relate distant parts of an
image to one another. A transformer takes the opposite stance: it treats the
image as a set of tokens and lets every token attend to every other through
self-attention~\cite{vaswani2017attention,dosovitskiy2021vit}, so global
relationships are available from the first layer at the cost of weaker built-in
spatial priors and a hunger for data. The \emph{attention} operation at the
heart of both transformers and our fusion modules computes a weighted average of
``value'' vectors, where the weights come from the similarity between a
``query'' and a set of ``keys''; when the query comes from one source and the
keys and values from another, the same mechanism becomes \emph{cross-attention},
which is how we inject clinical metadata into the image representation. Finally,
a Kolmogorov--Arnold Network (KAN)~\cite{liu2024kan} replaces the fixed
activation and linear weights of an ordinary layer with learnable spline
functions on the edges of the network, a recent idea we evaluate as a classifier
head.

\subsection{Review of Existing Research}

\subsubsection{Convolutional backbones used in this work}

The convolutional components of both our models are standard, well-understood
backbones, chosen because their behaviour and pretrained weights are reliable.
\textbf{MobileNetV2}~\cite{sandler2018mobilenetv2} introduced the inverted
residual block with a linear bottleneck, which keeps the network cheap by
expanding to a higher-dimensional space only inside each block and using
depthwise convolutions. Its early stages encode generic edges and textures very
efficiently, which is why we use a truncated MobileNetV2 as the stem of our
lightweight model. \textbf{EfficientNet-B0}~\cite{tan2019efficientnet} is the
base model of a family obtained by compound scaling of depth, width, and
resolution together; at roughly 5M parameters it is a strong accuracy-per-FLOP
baseline and one of the two backbones in our flagship model.
\textbf{DenseNet-121}~\cite{huang2017densenet} connects each layer to every
subsequent layer, encouraging feature reuse and giving a rich, gradient-friendly
representation; it is the second backbone in the flagship model and contributes a
different inductive bias from EfficientNet. We also include
\textbf{ResNet-18}~\cite{he2016resnet}, the residual-learning network that made
very deep training practical, as a familiar reference point.

\subsubsection{Transformers and CNN Transformer hybrids}

The \textbf{Vision Transformer}~\cite{dosovitskiy2021vit} showed that a pure
transformer, given enough data, can match convolutional models on image
classification by splitting an image into patches and treating them as a token
sequence. Its weakness on smaller datasets motivated hybrid designs that keep a
convolutional front end. \textbf{MobileViT}~\cite{mehta2022mobilevit} interleaves
convolutions with transformer blocks to get a mobile-friendly model that still
reasons globally, and \textbf{EfficientFormer}~\cite{li2022efficientformer} shows
that transformer-grade accuracy can be reached at MobileNet-like latency with a
carefully chosen dimension-consistent design. We treat both as direct,
same-spirit baselines, since a reviewer will reasonably ask how a new hybrid
compares against existing lightweight hybrids rather than against pure CNNs
alone. Our own design differs from these in where the convolution-to-transformer
hand-off happens and, more importantly, in that the transformer's class token is
used as a query that attends over clinical metadata.

\subsubsection{Kolmogorov Arnold Networks}

KANs~\cite{liu2024kan} are a recent alternative to the multilayer perceptron,
grounded in the Kolmogorov Arnold representation theorem, in which the learnable
parameters are the coefficients of B-spline functions placed on the network's
edges rather than scalar weights feeding fixed activations. The promise is more
expressive, more interpretable units, sometimes at fewer parameters. KANs are
still new and most evidence comes from small scientific-regression problems, so
whether they help as a classifier head on a noisy, imbalanced vision task is an
open question. We include a KAN head as the final classifier in our flagship
model to explore whether its learnable splines add value at the output stage of
a well-trained vision backbone.

\subsubsection{Clinical metadata fusion}

The work closest in spirit to ours is by Pacheco and
Krohling~\cite{pacheco2020impact,pacheco2021attention}. They showed first that
adding patient clinical information measurably improves automated skin cancer
detection~\cite{pacheco2020impact}, and then proposed an attention-based
mechanism to combine images and metadata inside a deep
model~\cite{pacheco2021attention}. Gessert et al.~\cite{gessert2020ensembles},
the basis of the top ISIC-2019 challenge entry, also fed metadata into an
ensemble of EfficientNets. The key methodological point we adopt from this line
of work is that missing metadata should be represented explicitly with a
learned ``missing'' embedding rather than imputed with a fabricated value,
so the model receives an honest ``unknown'' signal. Where we depart from prior
work is in the backbone: Pacheco and Krohling fused metadata onto heavy CNNs,
whereas we attach the same idea to a lightweight hybrid and to a dual-backbone
hybrid, and we make an explicit efficiency claim about the result.

\subsubsection{Skin lesion classification and ISIC benchmarks}

Beyond the specific components above, automated dermoscopy has a large
literature. Esteva et al.~\cite{esteva2017dermatologist} demonstrated
dermatologist-level melanoma classification with a single large CNN, and the
ISIC challenges~\cite{codella2018skin} have since standardised evaluation. The
ISIC-2019 dataset is itself a union of several sources, most prominently
HAM10000~\cite{tschandl2018ham10000} and BCN20000~\cite{isic2019challenge},
which is why source-aware splitting matters when building a fair evaluation. On
the ISIC-2019 leaderboard~\cite{isic2019leaderboard}, where the official ranking
metric is balanced accuracy, the strongest published entries are large
ensembles: the top entry (DAISYLab) reached a balanced accuracy of 0.636 with an
ensemble of multi-resolution EfficientNets plus SENet-154, followed by entries
around 0.61 and 0.59, again ensembles of EfficientNet and SE-ResNeXt variants.
These numbers set a useful ceiling and a reminder that single, lightweight
models operate in a different efficiency regime than the leaderboard's
ensembles.

Most directly relevant to our experimental design, Aktan et
al.~\cite{aktan2025comparison} recently compared several leading deep learning
architectures on this exact dataset, providing a clean, single-model reference
for how standard backbones behave on ISIC-2019 under a common protocol. Our
study extends that kind of controlled comparison in two directions: we add
clinical metadata fusion to the comparison, and we report results under a tight
efficiency budget, so that accuracy is always read together with parameter count
and compute.

